% Options for packages loaded elsewhere
\PassOptionsToPackage{unicode}{hyperref}
\PassOptionsToPackage{hyphens}{url}
%
\documentclass[
]{article}
\usepackage{amsmath,amssymb}
\usepackage{iftex}
\ifPDFTeX
  \usepackage[T1]{fontenc}
  \usepackage[utf8]{inputenc}
  \usepackage{textcomp} % provide euro and other symbols
\else % if luatex or xetex
  \usepackage{unicode-math} % this also loads fontspec
  \defaultfontfeatures{Scale=MatchLowercase}
  \defaultfontfeatures[\rmfamily]{Ligatures=TeX,Scale=1}
\fi
\usepackage{lmodern}
\ifPDFTeX\else
  % xetex/luatex font selection
\fi
% Use upquote if available, for straight quotes in verbatim environments
\IfFileExists{upquote.sty}{\usepackage{upquote}}{}
\IfFileExists{microtype.sty}{% use microtype if available
  \usepackage[]{microtype}
  \UseMicrotypeSet[protrusion]{basicmath} % disable protrusion for tt fonts
}{}
\makeatletter
\@ifundefined{KOMAClassName}{% if non-KOMA class
  \IfFileExists{parskip.sty}{%
    \usepackage{parskip}
  }{% else
    \setlength{\parindent}{0pt}
    \setlength{\parskip}{6pt plus 2pt minus 1pt}}
}{% if KOMA class
  \KOMAoptions{parskip=half}}
\makeatother
\usepackage{xcolor}
\usepackage{longtable,booktabs,array}
\usepackage{calc} % for calculating minipage widths
% Correct order of tables after \paragraph or \subparagraph
\usepackage{etoolbox}
\makeatletter
\patchcmd\longtable{\par}{\if@noskipsec\mbox{}\fi\par}{}{}
\makeatother
% Allow footnotes in longtable head/foot
\IfFileExists{footnotehyper.sty}{\usepackage{footnotehyper}}{\usepackage{footnote}}
\makesavenoteenv{longtable}
\setlength{\emergencystretch}{3em} % prevent overfull lines
\providecommand{\tightlist}{%
  \setlength{\itemsep}{0pt}\setlength{\parskip}{0pt}}
\setcounter{secnumdepth}{-\maxdimen} % remove section numbering
% unicode-fix.tex — render raw-Unicode symbols that the Latin Modern *text* font
% lacks, using Computer Modern math/text equivalents. This keeps the established
% Computer Modern look while filling the glyphs xelatex would otherwise drop.
% Included via: pandoc ... --include-in-header=unicode-fix.tex
\usepackage{newunicodechar}
\usepackage{pifont}
\usepackage{graphicx} % unicode-math's smallsetminus uses rotatebox
% Use the Computer Modern-compatible set-difference construction; the large
% Unicode glyph selected by TeX Live 2022 is absent from Latin Modern Math.
\AtBeginDocument{\let\setminus\smallsetminus}
% --- Greek letters -> math (Latin Modern Math; consistent with CM body text) ---
\newunicodechar{α}{\ensuremath{\alpha}}
\newunicodechar{β}{\ensuremath{\beta}}
\newunicodechar{γ}{\ensuremath{\gamma}}
\newunicodechar{δ}{\ensuremath{\delta}}
\newunicodechar{ε}{\ensuremath{\varepsilon}}
\newunicodechar{θ}{\ensuremath{\theta}}
\newunicodechar{μ}{\ensuremath{\mu}}
\newunicodechar{ν}{\ensuremath{\nu}}
\newunicodechar{ρ}{\ensuremath{\rho}}
\newunicodechar{σ}{\ensuremath{\sigma}}
\newunicodechar{τ}{\ensuremath{\tau}}
\newunicodechar{φ}{\ensuremath{\varphi}}
\newunicodechar{χ}{\ensuremath{\chi}}
\newunicodechar{ψ}{\ensuremath{\psi}}
\newunicodechar{η}{\ensuremath{\eta}}
\newunicodechar{κ}{\ensuremath{\kappa}}
\newunicodechar{λ}{\ensuremath{\lambda}}
\newunicodechar{π}{\ensuremath{\pi}}
\newunicodechar{ω}{\ensuremath{\omega}}
\newunicodechar{Ḣ}{\ensuremath{\dot H}}
\newunicodechar{≪}{\ensuremath{\ll}}
\newunicodechar{≫}{\ensuremath{\gg}}
\newunicodechar{⟺}{\ensuremath{\Longleftrightarrow}}
\newunicodechar{⟹}{\ensuremath{\Longrightarrow}}
\newunicodechar{⇔}{\ensuremath{\Leftrightarrow}}
\newunicodechar{⟨}{\ensuremath{\langle}}
\newunicodechar{⟩}{\ensuremath{\rangle}}
\newunicodechar{⁰}{\textsuperscript{0}}
\newunicodechar{¹}{\textsuperscript{1}}
\newunicodechar{²}{\textsuperscript{2}}
\newunicodechar{³}{\textsuperscript{3}}
\newunicodechar{⁴}{\textsuperscript{4}}
\newunicodechar{⁵}{\textsuperscript{5}}
\newunicodechar{⁶}{\textsuperscript{6}}
\newunicodechar{⁷}{\textsuperscript{7}}
\newunicodechar{⁸}{\textsuperscript{8}}
\newunicodechar{⁹}{\textsuperscript{9}}
% --- Relations / operators -> math ---
\newunicodechar{↔}{\ensuremath{\leftrightarrow}}
\newunicodechar{⇐}{\ensuremath{\Leftarrow}}
\newunicodechar{⇒}{\ensuremath{\Rightarrow}}
\newunicodechar{∂}{\ensuremath{\partial}}
\newunicodechar{∅}{\ensuremath{\varnothing}}
\newunicodechar{∈}{\ensuremath{\in}}
\newunicodechar{∼}{\ensuremath{\sim}}
\newunicodechar{≈}{\ensuremath{\approx}}
\newunicodechar{≠}{\ensuremath{\neq}}
\newunicodechar{≤}{\ensuremath{\leq}}
\newunicodechar{≥}{\ensuremath{\geq}}
\newunicodechar{⊂}{\ensuremath{\subset}}
\newunicodechar{⊃}{\ensuremath{\supset}}
% --- Special ---
\newunicodechar{ℏ}{\ensuremath{\hbar}}
\newunicodechar{ℤ}{\ensuremath{\mathbb{Z}}}
% --- Superscript modifier letters (footnote-style markers) -> text superscript ---
\newunicodechar{ʰ}{\textsuperscript{h}}
\newunicodechar{ʲ}{\textsuperscript{j}}
\newunicodechar{ᵃ}{\textsuperscript{a}}
\newunicodechar{ᵇ}{\textsuperscript{b}}
\newunicodechar{ᵈ}{\textsuperscript{d}}
\newunicodechar{ᵉ}{\textsuperscript{e}}
\newunicodechar{ᵍ}{\textsuperscript{g}}
\newunicodechar{ᶠ}{\textsuperscript{f}}
\newunicodechar{⁻}{\textsuperscript{\ensuremath{-}}}
% --- Subscript digits -> text subscript ---
\newunicodechar{₀}{\textsubscript{0}}
\newunicodechar{₁}{\textsubscript{1}}
\newunicodechar{₂}{\textsubscript{2}}
\newunicodechar{₃}{\textsubscript{3}}
\newunicodechar{₄}{\textsubscript{4}}
\newunicodechar{₈}{\textsubscript{8}}
% --- Checkmarks ---
\newunicodechar{✓}{\ding{51}}
\newunicodechar{✗}{\ding{55}}

% --- added 2026-08-02 after an independent glyph audit ---
\newunicodechar{ℓ}{\ensuremath{\ell}}
\newunicodechar{⊕}{\ensuremath{\oplus}}
\newunicodechar{⊗}{\ensuremath{\otimes}}
\newunicodechar{∠}{\ensuremath{\angle}}
\newunicodechar{≡}{\ensuremath{\equiv}}
\newunicodechar{∝}{\ensuremath{\propto}}
\newunicodechar{ℝ}{\ensuremath{\mathbb{R}}}
\newunicodechar{₅}{\ensuremath{_5}}
\newunicodechar{₋}{\ensuremath{_-}}
\newunicodechar{⁺}{\ensuremath{^+}}
\newunicodechar{⁽}{\ensuremath{^(}}
\newunicodechar{⁾}{\ensuremath{^)}}
% U+239B..U+23A0: multi-line parenthesis assembly pieces emitted by some
% converters. There is no sane inline equivalent; map to ordinary delimiters so
% the notation survives rather than vanishing silently.
\newunicodechar{⎛}{\ensuremath{\big(}}
\newunicodechar{⎜}{\ensuremath{\big\vert}}
\newunicodechar{⎝}{\ensuremath{\big(}}
\newunicodechar{⎞}{\ensuremath{\big)}}
\newunicodechar{⎟}{\ensuremath{\big\vert}}
\newunicodechar{⎠}{\ensuremath{\big)}}
\ifLuaTeX
  \usepackage{selnolig}  % disable illegal ligatures
\fi
\IfFileExists{bookmark.sty}{\usepackage{bookmark}}{\usepackage{hyperref}}
\IfFileExists{xurl.sty}{\usepackage{xurl}}{} % add URL line breaks if available
\urlstyle{same}
\hypersetup{
  pdftitle={Observational Incompleteness: From Deterministic Dynamics to Quantum Observations},
  pdfauthor={Alex Maybaum},
  hidelinks,
  pdfcreator={LaTeX via pandoc}}

\title{Observational Incompleteness: From Deterministic Dynamics to
Quantum Observations}
\author{Alex Maybaum}
\date{October 2026}

\begin{document}
\maketitle

\textbf{Status:} DRAFT PRE-PRINT\\
\textbf{Classification:} Theoretical Physics / Foundations

\begin{center}\rule{0.5\linewidth}{0.5pt}\end{center}

\hypertarget{abstract}{%
\subsection{Abstract}\label{abstract}}

Observational incompleteness (OI) is the hypothesis that every physical
observer is a proper part of a finite reversible deterministic system
and records only the visible marginals of its dynamics. At every finite
horizon, three classes of law on visible trajectories coincide: all
finite stochastic laws, the visible marginals of finite reversible
dynamics, and the fixed-basis Born records of finite unitary systems.
The components of this correspondence are known, from the
trajectory-labelling realization of Egri et al.~and Barandes's treatment
of permutation dynamics as unistochastic processes; it is stated here as
one theorem. Each class contains the Markov laws, so representability
does not distinguish a law with memory from one without, and the Born
exponent is not selected. The discriminating content concerns memory:
every deterministic completion carries the visible memory in its hidden
state, and on a fixed finite reversible system finite recurrence turns
history readback into stochastic indivisibility somewhere in the
recurrence cycle (each with one step not machine-checked). Exact finite
operational quantum mechanics is characterized by five completion
conditions, none implied by OI alone. Open are which of conditions
(iii)--(v) a physical substratum supplies, the selection of the
evolution relating two times, and any test of the hypothesis. Applied to
the cosmological constant, OI places the zero-point sum and the
effective cosmological constant at different levels of description;
whether the zero-point sum sources curvature depends on covariant field
equations that are not derived, so this is a proposal, not a resolution.
Applied to black-hole information, finite computations separate the
conservation of information under bijective dynamics from its recovery
by an observer with restricted access; the framework's Page curve is not
established, and the information problem is not resolved.

\hypertarget{introduction-observational-incompleteness-and-the-emergence-of-quantum-behaviour}{%
\subsection{1. Introduction: Observational Incompleteness and the
emergence of quantum
behaviour}\label{introduction-observational-incompleteness-and-the-emergence-of-quantum-behaviour}}

\hypertarget{embedded-observation}{%
\subsubsection{1.1 Embedded observation}\label{embedded-observation}}

Embedded observation is modelled by a triple \((S,\varphi,V)\) {[}Main
§1.2{]}: a finite total system \(S=\mathcal C_V\times\mathcal C_H\) of
observer configurations \(\mathcal C_V\) and complement configurations
\(\mathcal C_H\), a bijection \(\varphi\) of \(S\), and an observer
\(V\) that registers \(x\in\mathcal C_V\) and not \(h\in\mathcal C_H\).

The bijection \(\varphi\) is the reversible representative of the hidden
dynamics, motivated by observed near-unitarity, and finiteness of all of
\(S\) is posited {[}Main §1.2, Remark (status of the early
commitments){]}. Hidden states compatible with the same \(x\) send it to
different futures, so the visible sector evolves stochastically although
the whole evolves deterministically.

Observational incompleteness (OI) is the physical hypothesis that every
physical observer is in this situation. Identifying nature with such a
system is a representation choice, motivated empirically; the theorems
below are mathematical, and no test of the identification at the level
of memory witnesses is established {[}Main §4.5{]}. The incompleteness
is inaccessibility of the hidden sector, not joint unsharpness of
non-commuting observables: it is present for a single fixed readout, and
no uncertainty relation constrains fixed-readout visible laws {[}Main
§3.4{]}.

\hypertarget{the-idea-and-its-two-consequences}{%
\subsubsection{1.2 The idea and its two
consequences}\label{the-idea-and-its-two-consequences}}

Marginalizing a reversible dynamics over an unobserved register produces
stochastic laws, conditional memory, and failures of stochastic
divisibility, in which no stochastic intermediate transition map exists
(where the candidate intermediate matrix exists, it has negative
entries). The stochastic--quantum correspondence relates these features
to quantum theory {[}Bar23; Bar25{]}. At the level of finite-horizon
observable laws the correspondence of §2 (Theorem 1) is exact and
universal, so the content lies in where the memory of a law must sit
(§2.4) and what finite recurrence does with it (§2.5). The paper
develops one idea, OI: its foundation and quantum-emergence results,
which stop short of operational quantum mechanics (§§2--3), and two of
its consequences, for the cosmological constant (§4) and for black-hole
information (§5). Neither application is a demonstrated resolution of
its problem; §6 states what would decide each. Proofs, the status ledger
and the probe code are in {[}Main{]}, which this paper introduces;
gravitational applications are in {[}GR{]} and {[}Book{]}.

\textbf{Contributions.} The paper claims no new theorem or physical
result; its contribution is expository. It presents the established
mathematical results of OI together, separates the mathematical
correspondence from its physical interpretation, and states two
consequences with their conditions and limitations. The black-hole
results of §5 are reported from unpublished repository research records
{[}BHa; BHb{]}, with written proofs and exact or exhaustive checks,
whose novelty has not been established. The components of Theorem 1 are
known {[}EGGH26; Bar23; Bar25; Bar26; SGD19{]}; the paper states them as
one theorem with each direction attributed. The memory bounds of §2.4
are forms of the data-processing and Pinsker inequalities. The
recurrence argument of §2.5 combines finite order with the known
contraction criterion for divisibility {[}RHP14; BD16{]}, stated
concurrently by Doukas {[}Dou26{]}. The results of §3.2--§3.3 are proved
in {[}GR{]} and {[}Main{]}.

\emph{Conventions.} ``Machine-checked'' means checked in Lean 4/Mathlib
{[}dMU21; Mathlib20{]} at commit \texttt{17b830bc} by the recipe under
``Code, data and formal verification'' (declarations in Appendix A);
``proved on paper'' means a written proof without a machine check; an
``exact probe'' is an exhaustive computation in exact rational
arithmetic at small size. A \emph{faithful realization} of a law is a
bijection \(\varphi\) of \(\mathcal C_V\times\mathcal C_H\) with an
initial law under which the visible trajectory has that law; the finite
reversible system \((S,\varphi)\) is the \emph{substratum}.

\hypertarget{theoretical-foundation-the-established-results-and-their-assumptions}{%
\subsection{2. Theoretical foundation: the established results and their
assumptions}\label{theoretical-foundation-the-established-results-and-their-assumptions}}

\hypertarget{setting}{%
\subsubsection{2.1 Setting}\label{setting}}

Fix finite sets \(\mathcal C_V\) (with \(|\mathcal C_V|\ge2\)) and
\(\mathcal C_H\), a bijection \(\varphi\) of
\(\mathcal C_V\times\mathcal C_H\), and a horizon \(K\). A law is a
probability distribution on visible trajectories
\((x_0,\dots,x_K)\in\mathcal C_V^{K+1}\).

The visible law of a realization depends on the dynamics, the partition
and the initial law; the prior is part of the realization datum, not
selected by invariance under \(\varphi\), and the counting measure is a
maximal-entropy default {[}Main §1.2, Remark (the measure as a
realization datum){]}.

Three classes of law are compared {[}Main §3.4{]}:

\begin{itemize}
\tightlist
\item
  \textbf{\(\mathsf S_K\)}: all probability distributions on
  \(\mathcal C_V^{K+1}\).
\item
  \textbf{\(\mathsf D_K\)}: the visible trajectory laws of a bijection
  of \(\mathcal C_V\times\mathcal C_H\), for some finite
  \(\mathcal C_H\), under an initial law that may correlate the visible
  and hidden coordinates.
\item
  \textbf{\(\mathsf Q_K\)}: the laws of finite-dimensional unitary
  systems read out by projective Born measurement in one fixed basis at
  every step, from an initial state diagonal in that basis, the record
  being a fixed function of each outcome.
\end{itemize}

\hypertarget{the-finite-horizon-correspondence}{%
\subsubsection{2.2 The finite-horizon
correspondence}\label{the-finite-horizon-correspondence}}

\textbf{Theorem 1} (finite-horizon stochastic--reversible--unitary
equivalence {[}Main §3.4{]}). \emph{For every finite visible alphabet
and every finite horizon \(K\),} \[
\mathsf S_K=\mathsf D_K=\mathsf Q_K ,
\] \emph{so each class is the set of all probability laws on
\(\mathcal C_V^{K+1}\).}

Theorem 1 is machine-checked as \texttt{finite\_horizon\_equivalence},
with the inclusions \(\mathsf S_K\subseteq\mathsf D_K\)
(\texttt{S\_imp\_D}), \(\mathsf D_K\subseteq\mathsf Q_K\)
(\texttt{D\_imp\_Qfb}) and \(\mathsf Q_K\subseteq\mathsf S_K\)
(\texttt{Qfb\_imp\_S}), and the reverse inclusions as their
compositions; no construction enters as a hypothesis.

\textbf{Realization.} For \(\mathsf S_K\subseteq\mathsf D_K\) the
realization stores a sampled trajectory, with a clock, in the hidden
register {[}Main §3.4, Theorem (finite-horizon
stochastic--reversible--unitary equivalence); EGGH26, Prop. 13{]}. For a
law with rational probabilities, a realization whose hidden prior is
independent of \(x_0\) also exists {[}Main §3.4, Theorem (finite-horizon
process dilation){]}; it is proved on paper, with exact probes.

\textbf{Quantum representation.} For \(\mathsf D_K\subseteq\mathsf Q_K\)
the representation is the permutation unitary
\(U_\varphi|x,h\rangle=|\varphi(x,h)\rangle\), read out in the
configuration basis at every step {[}Main §3.4{]};
\(\mathsf Q_K\subseteq\mathsf S_K\) is the Born rule.

\textbf{Attribution.} The trajectory-labelling realization is that of
Egri et al.~{[}EGGH26, Prop. 13{]}, packaged as one iterated bijection;
the independent-prior form is their Prop. 16, without reversibility, and
its Markov one-step case is the random-mapping representation {[}LP17,
Prop. 1.5{]}. A permutation dynamics is a special case of a
unistochastic process {[}Bar25, p.~12; Bar23 §4.3; Bar26 §4{]}, and the
one-step \(n^2\)-dimensional ancilla dilation of {[}Main §3.1{]}
coincides with the constructions of Barandes {[}Bar26, Thm 3.3{]} and
Schmidt {[}Sch21, Thm 2{]}. The every-step projectively measured record
as a probability law is in {[}SGD19, Eq. (3){]}; that work and
{[}Milz20{]} study whether such records are consistent with omitting
measurements, a question \(\mathsf Q_K\), which measures at every step,
does not meet.

\hypertarget{scope-of-the-correspondence}{%
\subsubsection{2.3 Scope of the
correspondence}\label{scope-of-the-correspondence}}

Markov laws lie in all three classes. So neither a finite reversible
realization nor a fixed-basis unitary representation distinguishes a law
with memory from one without {[}Main §3.4, Remark (the conditions are
diagnostics, not hypotheses){]}.

The relation \(p=|U|^2\) is part of the definition of \(\mathsf Q_K\),
not an output. The unitaries supplied are permutations, with entries 0
and 1, so the recovered table is the same for \(|U|^p\) at every
\(p>0\). The quadratic exponent is therefore not selected {[}Main §3.4;
Bar26, Thm 3.1{]}.

A finite bijection has finite order \(L\), so one fixed realization has
an \(L\)-periodic multi-time law: the correspondence holds per horizon,
and a nonperiodic law has no exact all-time realization by one fixed
finite substratum {[}Main §3.4{]}.

\hypertarget{where-the-memory-lives}{%
\subsubsection{2.4 Where the memory
lives}\label{where-the-memory-lives}}

Write \(M_t := I(X_{<t};X_{t+1}\mid X_t)\) for the conditional memory at
time \(t\). In a realization, write \(f_x(h)=\pi_V(\varphi(x,h))\), with
\(\pi_V\) the projection onto \(\mathcal C_V\), and \(\mu_p\) for the
hidden posterior after a visible history \(p\) ending at \(x\).

\textbf{Theorem 2} (unavoidable hidden predictive memory {[}Main
§3.4{]}). \emph{In every faithful realization:}

\begin{itemize}
\tightlist
\item
  \emph{(a) \(P(X_{t+1}\mid p,x)=(f_x)_\#\mu_p\);}
\item
  \emph{(b) for histories \(p,p'\) ending at \(x\),
  \(\|P(X_{t+1}\mid p,x)-P(X_{t+1}\mid p',x)\|_{\mathrm{TV}}\le\|\mu_p-\mu_{p'}\|_{\mathrm{TV}}\);}
\item
  \emph{(c) \(M_t\le I(X_{<t};H_t\mid X_t)\le\log_2|\mathcal C_H|\).}
\end{itemize}

(a)--(c) are machine-checked for a joint law of history, present and
hidden state. Their identification with a realization's time-\(t\) law
is not machine-checked. The written proof is in {[}Main §3.4, Theorem
(unavoidable hidden predictive memory){]}; see also {[}CT06{]}.

\textbf{Corollary} (proved on paper {[}Main §1.5{]}). \(M_t>0\) holds
exactly when every faithful realization carries readout-relevant hidden
memory at \(t\), that is, a pair of posteriors with different
pushforwards. The forward direction rests on Theorem 2(a) in each
realization; the converse on Theorem 2(a) in one realization, which
Theorem 1 supplies (\texttt{S\_imp\_D}); both use the standard vanishing
criterion for conditional mutual information {[}CT06{]}, which, like the
time-\(t\) identification, is not machine-checked.

Two further bounds tie memory to the hidden sector. A realization of a
law whose conditional past--future information reaches \(I^*\) needs
\(|\mathcal C_H|\ge2^{I^*}\) (machine-checked in the same form as
Theorem 2). If two values of \(X_{<k}\) followed by a common present
state \(X_k\), each of joint probability at least \(p_0\), induce
next-step laws at total-variation distance at least \(\delta\), then
\(I(X_{<k};X_{k+1}\mid X_k)\ge p_0\delta^2/\ln2\) (in bits), by
Pinsker's inequality (proved on paper, with exact probes) {[}Main §2.3;
CT06{]}.

{[}Main §1.3{]} names four conditions on a hidden sector: (C1)
influence, (C2) persistence of records, (C3) capacity, and (C4) history
readback, in which two histories with the same present induce different
next-step laws. In a faithful realization (C4) coincides with \(M_t>0\)
by the vanishing criterion {[}CT06{]}, and the Corollary places it in
the hidden state; (C1) and (C3) follow from it, and (C2) is a
physical-regime premise. The conditions are diagnostics of a
realization, not hypotheses of Theorem 1.

\hypertarget{recurrence-and-indivisibility}{%
\subsubsection{2.5 Recurrence and
indivisibility}\label{recurrence-and-indivisibility}}

Write \(\Gamma_t\) for the \emph{rooted} visible transition matrix from
time 0 to time \(t\) (rows indexed by \(x_0\)). The family is
\emph{P-divisible on \(\{0,\dots,n\}\)} if for all \(s<t\le n\) some
stochastic \(\Lambda\) satisfies \(\Gamma_t=\Gamma_s\Lambda\). A
stochastic divisor contracts total variation, so under P-divisibility
\(\|\delta_a\Gamma_t-\delta_b\Gamma_t\|_{\mathrm{TV}}\) cannot increase.
This is the divisibility criterion of {[}RHP14, Thm 2.4{]} and the
data-processing direction of {[}BD16, Prop. 1{]}.

\textbf{Theorem 3} (recurrence and overlap {[}Main §2.3{]}). \emph{Let
\(\varphi\) be a bijection of \(\mathcal C_V\times\mathcal C_H\) with a
fixed hidden prior. If two rows of \(\Gamma_s\) overlap in a column and
\(\Gamma_n=I\) for some \(n>s\), the family is not P-divisible on
\(\{0,\dots,n\}\).}

Theorem 3 is machine-checked in its rooted form, and Doukas states the
same overlap-and-return argument concurrently {[}Dou26, Prop. 1{]}.
Monotone distance is necessary for P-divisibility but not sufficient
{[}Main §2.3{]}.

A finite square stochastic matrix with a stochastic one-sided inverse is
a permutation matrix {[}BB12, Problem 23{]} (machine-checked). With this
lemma, on a fixed finite reversible substratum, history readback forces
indivisibility somewhere in the recurrence cycle {[}Main §2.3, Theorem
(history readback + finite recurrence implies indivisibility){]}. The
conclusion is machine-checked, by the route of Theorem 3, from a
routed-readback form of (C4), in which a visible-history record written
into the hidden state is read back later {[}Main §1.3, the realization
clause of (C4){]}; from history-level (C4) it is proved on paper.

Appendix B works examples: a minimal model of Theorem 3 with six hidden
states, memory without short-window indivisibility, and a two-card game,
the smallest realization of Theorem 3.

\hypertarget{quantum-emergence-what-is-demonstrated-and-what-remains-open-for-equivalence-with-quantum-mechanics}{%
\subsection{3. Quantum emergence: what is demonstrated, and what remains
open for equivalence with quantum
mechanics}\label{quantum-emergence-what-is-demonstrated-and-what-remains-open-for-equivalence-with-quantum-mechanics}}

\hypertarget{what-is-demonstrated}{%
\subsubsection{3.1 What is demonstrated}\label{what-is-demonstrated}}

The indivisibility of §2.5 is global: it occurs somewhere in the
recurrence cycle, not on every window. Conditional memory does not force
short-window indivisibility (Appendix B).

Three notions are distinct, and no unconditional implication holds among
them {[}Main §3.2, Remark (three notions, kept apart){]}: (1) accessible
memory, \(M_t>0\); (2) stochastic P-indivisibility of the visible
population law; (3) quantum non-Markovianity under a chosen criterion.
Tracing out a permutation unitary from a diagonal hidden state gives
channels that preserve diagonal states, and for them (2) implies
CP-indivisibility {[}Main §3.2, Theorem (CP-indivisibility, framework
scope){]}. Outside that class the implication fails: with \(\Phi_{t_1}\)
conjugation by the Hadamard gate and \(\Phi_{t_2}=\mathrm{id}\) the
family is CP-divisible, while its computational-basis populations,
\(\tfrac12J\to I\) with \(J\) the all-ones matrix, are P-indivisible
{[}Main §3.2, Remark (scope: P-indivisibility is not quantum
non-Markovianity){]}. Distinguishability backflow witnesses divisibility
failure in one direction only {[}BLP09{]}; for bijective dynamical maps
on a \(d\)-dimensional system, CP-divisibility is equivalent to monotone
distinguishability with a \((d+1)\)-dimensional ancilla {[}BJA17{]}.

Theorem 1 is a representation result: it states which laws admit
realizations and quantum representations, not why a physical system
exhibits a law with memory {[}Main §3.4, Remark (representation, not
explanation){]}. Calvo states a proof of the six axioms of textbook
quantum mechanics from the single axiom that every physical system
evolves according to a generally indivisible stochastic law {[}Cal26{]};
Theorem 1 is a correspondence between classes of law and derives no such
axioms.

\hypertarget{from-the-correspondence-to-operational-quantum-mechanics}{%
\subsubsection{3.2 From the correspondence to operational quantum
mechanics}\label{from-the-correspondence-to-operational-quantum-mechanics}}

Operational quantum mechanics requires, beyond Theorem 1's fixed-basis
statistics, coherent preparations, noncommuting interventions and
composites represented by one instrument theory.

Fix a nonempty finite system and an operational theory specifying which
finite outcome families of linear maps are available on the system and
at every ancilla level. The theory contains the identity and is closed
under coarse-graining of outcomes and, at every ancilla level, under
outcome-conditioned composition; it has an available preparation of a
uniformly mixed ancilla, which an available operation may follow, an
available spectator-independent local readout of the ancilla, and
ancilla discard: an available preparation followed by an available
instrument, with the ancilla then discarded, is an available instrument
on the system with the same outcomes {[}GR §3.3{]}. The five conditions
are: (i) available families send positive states to positive states and
their outcome traces sum to the input trace; (ii) operations remain
available with a one-state ancilla adjoined; (iii) an untouched
independent system can be adjoined without changing an intervention;
(iv) every finite unitary on every composite is available; (v) attaching
a fresh uniformly prepared ancilla to a composite, acting, and
discarding it gives an available operation.

\textbf{Theorem 4} (operational-completion characterization {[}GR
§3.3{]}). \emph{For every nonempty finite system, the available outcome
families on the system and at every positive ancilla level are exactly
the normalized finite Kraus instruments if and only if conditions
(i)--(v) hold.}

Both directions are machine-checked, \(\Rightarrow\) as
\texttt{physical\_of\_exactAll} and \(\Leftarrow\) as
\texttt{exactAll\_of\_physical\_general}, for every nonempty finite
carrier.

On the qubit, each condition is independent of the other four and of a
fixed finite observation process satisfying (C1)--(C4). For each
condition there is a qubit theory that realizes that process and the
other four conditions, and fails it. In particular, a qubit theory that
never creates coherence realizes the observation process and is not
quantum. So OI alone does not imply quantum mechanics. Which of
(iii)--(v) a physical substratum supplies is open {[}GR §3.3{]}.

\hypertarget{the-evolution-relating-two-times-is-not-selected}{%
\subsubsection{3.3 The evolution relating two times is not
selected}\label{the-evolution-relating-two-times-is-not-selected}}

Representability does not fix the unitary evolution relating two times
{[}Main §3.4, Remark (scope: representability, not selection of the
relating evolution){]}. Two coherent lifts of one visible family,
unitary dilations that reproduce the visible law at every time and
compose exactly across times, can differ by a time-dependent right
action invisible to the visible law and then have different relating
evolutions, while a constant deformation cancels from every relating
evolution (both machine-checked).

Take a lift with trivial ancilla at one time, and form the fixed-basis
datum whose evolution is the lift's unitary, with uniform initial law
and identity readout. If two such data have the same one-step visible
matrix, their root-conditioned trajectory laws agree at every horizon
(\texttt{bridge\_traj}); conversely, if their rooted families agree,
their one-step visible matrices are equal
(\texttt{rooted\_eq\_iff\_slice\_eq}, forward direction). Each such
trajectory law is fixed-basis realizable and stochastic
(\texttt{realData\_traj\_stochastic}). The fixed-basis correspondence
therefore acts on the quotient by visible equality. A selection
condition beyond representability and ordinary coherence would be
needed, and none is supplied here.

\hypertarget{no-fixed-finite-carrier-reproduces-a-qubit}{%
\subsubsection{3.4 No fixed finite carrier reproduces a
qubit}\label{no-fixed-finite-carrier-reproduces-a-qubit}}

A fixed finite realization with \(N\) ontic configurations generates
only behaviour matrices that factor nonnegatively through \(N\) states.
Simulating a qubit's prepare-and-measure statistics needs classical
memory unbounded over finite experiments: explicit tasks are given by
Galvão and Hardy {[}GH03{]} and Montina {[}Mon08{]}, and Montina reports
Hardy's result that the number of ontic states cannot be finite
{[}Mon08; Har04{]}. So no fixed finite realization reproduces the full
prepare-and-measure theory of a qubit exactly; the factorization is
proved, and the unboundedness is imported {[}Main §3.4, Proposition
(finite-substratum operational obstruction --- classical-dimension
form); HKLP24, Thm 4{]}. The obstruction does not touch Theorem 1, whose
realization is chosen per law.

\hypertarget{the-born-exponent-and-bell-structure}{%
\subsubsection{3.5 The Born exponent and Bell
structure}\label{the-born-exponent-and-bell-structure}}

The origin of the quadratic Born exponent, which the correspondence does
not select (§2.3), is open {[}Main §3.4; Bar26, Thm 3.1{]}.

A deterministic completion that keeps pointwise causal cones and a
setting-independent pre-setting ensemble is Bell-local,
\(|S_{\mathrm{CHSH}}|\le2\) {[}Bel64; CHSH69; Main §3.3{]}, so a
deterministic completion reproducing quantum Bell correlations must give
up one of the two. The framework keeps measurement independence and
adopts ontic parameter dependence, in which one wing's microstate
response depends on the remote setting without operational signalling,
and the compatibility of that choice with operational no-signalling and
the gravitational continuum limit of the substratum's spatial coupling
graph is the open hypothesis H-Bell {[}Main §3.3{]}.

\hypertarget{the-cosmological-constant-how-oi-might-change-the-interpretation-of-vacuum-energy}{%
\subsection{4. The cosmological constant: how OI might change the
interpretation of vacuum
energy}\label{the-cosmological-constant-how-oi-might-change-the-interpretation-of-vacuum-energy}}

Summing zero-point energies of the quantum fields up to a Planck-scale
cutoff gives a vacuum energy density some 120 orders of magnitude above
the observational value {[}Wei89, §I; §III, Eqs. (3.4)--(3.6){]} (the
mismatch), the exact power depending on the cutoff convention; if that
density sources curvature, it must be cancelled to that accuracy
{[}Wei89, §III{]}.

In OI (§1.1) the zero-point sum of the emergent quantum description
(§§2--3) and the effective cosmological constant are defined at
different levels of description. A uniform shift of the emergent
Hamiltonian changes no visible transition probability, so the offset has
no content in the visible statistics of the quantum description {[}GR
§6.2{]}. The invariance is the familiar statement that no
non-gravitational measurement sees the absolute vacuum energy; OI's
proposal is that the mismatch requires no cancellation, and its content
lies in how the two levels couple to geometry.

Whether the offset also fails to source curvature is a separate coupling
question. Settling it requires covariant field equations derived from
the substratum, with their coupling to sources (Level G3 of {[}GR{]}),
which are open {[}GR §6.2{]}. The quantum description therefore does not
decide between gravitational completions that differ by a term
\(\rho_C\,g_{\mu\nu}\), such as
\(\mathcal G_{\mu\nu}=8\pi G\,T^{\rm cl}_{\mu\nu}\) and
\(\mathcal G_{\mu\nu}=8\pi G\,(T^{\rm cl}_{\mu\nu}-\rho_C\,g_{\mu\nu})\),
with \(\mathcal G_{\mu\nu}\) the gravitational side of the field
equations (the Einstein tensor in general relativity),
\(T^{\rm cl}_{\mu\nu}\) the classical stress-energy of the substratum
and \(\rho_C\) the energy density of the offset {[}GR §6.2{]}. The
proposal that the mismatch requires no cancellation is therefore
conditional on that open step.

If the covariant field equations take the form {[}GR{]} targets,
\(E_{\mu\nu}[g,\mathrm{OI}]=8\pi G\,(T_{\mu\nu}-\mathcal F[g,T]\,g_{\mu\nu})\),
with \(E_{\mu\nu}[g,\mathrm{OI}]\) the gravitational side still to be
derived, \(\mathcal F[g,T]\) a scalar functional of the metric and the
stress-energy, and the classical stress-energy of the substratum as
source, the zero-point sum of the emergent description is not a source
and needs no cancellation; the vacuum term that remains, from the
substratum's classical energy density and from \(\mathcal F\), is not
computed. It does not explain the observed value of Λ, which enters as a
cosmological input. The scale \(\rho\sim H_{\rm Hub}^2/G\), with
\(H_{\rm Hub}\) the Hubble rate, is the Friedmann normalization of the
total density, and the fraction carried by Λ is an input.

With a Planck-scale cutoff the mismatch equals the de Sitter entropy up
to a numerical factor fixed by the cutoff convention; this is a
dimensional identity, not an independent prediction. How the energy of
electroweak and QCD vacuum transitions couples to gravity is not fixed
by the descriptive distinction and requires the same open coupling step.

Removing the vacuum energy of matter as a source of curvature is also
the aim of the trace-free Einstein equations, which leave the observed
value of Λ unexplained {[}EvEMU11{]}, and of vacuum-energy sequestering,
which also renders contributions from phase transitions small
{[}KP14{]}; Volovik argues from the thermodynamics of quantum many-body
ground states that the ultraviolet zero-point sum does not set the
vacuum energy {[}Vol06{]}. The present proposal differs in locating the
gravitational source in the substratum microstate, and its coupling step
is open.

A derivation of the covariant field equations, with their conservation
and source coupling, from the substratum would turn the proposal into a
theorem or refute it. Until then the proposal reinterprets the problem
and does not resolve it.

\hypertarget{black-hole-information-preservation-versus-accessible-recovery}{%
\subsection{5. Black-hole information: preservation versus accessible
recovery}\label{black-hole-information-preservation-versus-accessible-recovery}}

The black-hole information problem asks whether the radiation of an
evaporating black hole carries the information of what formed it
{[}Haw76; Pag93; Pag13{]}. Because the dynamics of OI is a bijection
observed through a partition (§1.1), two questions separate: whether
information is conserved in the global reversible dynamics, and whether
it can be recovered from the records an observer can access, its
\emph{access class}. The results below are finite computations in toy
models of configurations, permutations and configuration reads. They
concern the framework's nested-partition construction {[}GR App. A; Book
§7.6{]}, not a particular black hole.

Under OI, information is conserved in the global reversible dynamics
because the dynamics is a bijection, which is a premise of the
framework. Conservation neither implies nor is necessary for recovery
from the records an observer can access: finite exhaustive computations
exhibit bijections with no accessible advantage, that is, no gain in
success probability over a blind guess, and non-injective maps with
exact recovery {[}BHa{]}.

Let \(\mathcal U\) be the set of joint values of the registers left
unread at the end and \(\mathcal M\) the message set. Of the 40,320
bijections of \(\{0,1\}^3\), 1,376 give zero accessible advantage for a
binary message with one register unread after two emissions. For
bijective dynamics, complete endpoint access gives exact recovery, and
with \(\mathcal U\) inaccessible \(P_{\mathrm{succ}}\ge1/|\mathcal U|\),
so zero advantage in a bijective instance requires
\(|\mathcal U|\ge|\mathcal M|\); 32 non-injective maps of \(\{0,1\}^2\)
violate the bound on \(P_{\mathrm{succ}}\) {[}BHa{]}. These endpoint
guarantees are sufficient, not necessary. The bound is the
remnant-capacity form of a known fact, and its linear case is the
secret-sharing threshold {[}BHa; Sha79{]}.

The Page-curve calculation of {[}GR App. A{]} addresses recovery under
its own conditions and is a conditional result {[}Book §7.6, §18.5{]}.
Its configuration space is
\(\mathcal C_V\times\mathcal C_B\times\mathcal C_D\), with \(V\) the
exterior visible region, \(B\) the black-hole interior and \(D\) the
sector beyond the cosmological horizon; the radiation \(R\), a subsystem
of \(V\), is emitted from \(B\) {[}GR App. A.1, A.7{]}. For all but an
exponentially negligible fraction of initial conditions, the full-cycle
time average of the radiation's entanglement entropy, at the bipartition
of each stage held fixed, equals the Page value
\(S_{\rm Page}(d_B,d_R)\) for factor dimensions \(d_B,d_R\), up to the
correction term of Theorem A.8 {[}Pag93; GR, Thms A.8--A.9{]}. This
holds under the scrambling hypothesis H-scramble, that the orbit
ensemble equidistributes over quantum pure states on the shell;
H-scramble fails for every bijection acting on a definite initial
configuration with the bipartition aligned to the configuration
coordinates and is not established {[}GR App. A.7; BHb, C1,
C6(a)--(b){]}. It also needs the decoupling condition H-decouple:
\(B\cup R\) evolves with no coupling to \(D\) or to the rest of the
visible sector, \(V\setminus R\), and its joint state is pure, over
whole cycles of the dynamics; timescale ordering does not supply it
{[}BHa, item F{]}, and it is not established {[}GR App. A.7{]}. Theorem
A.8 of {[}GR{]} takes the whole configuration space
\(\mathcal C_B\times\mathcal C_R\) as a single energy shell. That
requires the energy to be constant there, so that energy conservation
places no constraint on the orbit, and it does not follow from the
conservation of a Hamiltonian that takes more than one value there. An
evaporating black hole whose radiation carries away its energy is in the
proper-shell case, where the mean entropy depends on the shell and in
general differs from the Page value; that case is open {[}GR App. A.7;
BHb, C6(c){]}. The result is a time average at a fixed bipartition, and
whether the entropy along a single trajectory follows the same curve is
open {[}GR App. A.7; BHb, C6(d){]}. In the dependency ladder of {[}GR{]}
the Page curve is also conditional on the covariant field equations. The
generalized second law is not derived {[}GR App. A.6{]}; agreement with
quantum-extremal-surface computations {[}Pen20{]} is not established;
the correspondence between islands and substratum cycles is conjectural
{[}Book §7.6{]}. The Page curve of {[}GR App. A{]} is therefore not
established, and the black-hole information problem is not resolved
here.

\hypertarget{discussion-and-conclusion-implications-limitations-and-decisive-open-tests}{%
\subsection{6. Discussion and conclusion: implications, limitations and
decisive open
tests}\label{discussion-and-conclusion-implications-limitations-and-decisive-open-tests}}

\hypertarget{implications}{%
\subsubsection{6.1 Implications}\label{implications}}

If the hypothesis holds, the fixed-basis statistics of quantum theory
are the visible laws of an embedded observer, and their departures from
Markov behaviour are memory carried in an inaccessible sector; whether
the rest of operational quantum mechanics follows depends on completion
conditions that OI does not imply (§3.2). Quantum Bell correlations are
not accounted for by memory alone: with measurement independence kept,
they need ontic parameter dependence, whose compatibility with
operational no-signalling and the continuum limit is the open hypothesis
H-Bell (§3.5).

\hypertarget{what-is-not-claimed}{%
\subsubsection{6.2 What is not claimed}\label{what-is-not-claimed}}

\begin{enumerate}
\def\labelenumi{\arabic{enumi}.}
\tightlist
\item
  No derivation of quantum mechanics: Theorem 1 contains every finite
  law, and the Born exponent is not selected.
\item
  No unique selection of the operational theory: conditions (i)--(v) are
  not implied by OI alone, their independence is shown on the qubit,
  which of (iii)--(v) a physical substratum supplies is open, and the
  visible law does not fix the relating evolution.
\item
  No characterization of memory by the existence of a realization
  (§2.3), no identification of P-indivisibility with quantum
  non-Markovianity, and no claim that memory forces indivisibility on
  short windows (§3.1).
\item
  No resolution of the black-hole information problem: the Page curve of
  {[}GR App. A{]} is not established (§5), and the results of §5
  reported from {[}BHa; BHb{]} concern finite toy models, each proved on
  paper or checked exhaustively at small size or by exact arithmetic. No
  claim is made that information reaches \(D\), that {[}GR, Thm A.7{]}
  fails within its stated scope, or that no other decoupling argument
  exists.
\item
  No claim that the physical hypothesis is established or tested.
\item
  No claim that the card game (Appendix B) produces quantum rather than
  classical statistics: on card inputs its statistics are classical
  stochastic ones; under the unitary embedding its one-step channel is
  entanglement-breaking and its two-step family admits a
  classical-memory decomposition; it violates no Bell inequality.
\item
  No prediction of the value of Λ or \(\Omega_\Lambda\), and no
  resolution of the cosmological constant problem: §4 is conditional on
  the covariant field equations.
\item
  No claim that the viewpoint is new: deterministic systems analysed
  with quantum methods have been studied before, for example in
  {[}tH14{]}. What the paper states are the cited theorems and their
  scopes.
\end{enumerate}

\hypertarget{decisive-open-tests}{%
\subsubsection{6.3 Decisive open tests}\label{decisive-open-tests}}

For quantum emergence: whether a physical substratum supplies completion
conditions (iii)--(v), and whether history readback (C4) holds at the
cosmological partition {[}GR §2.2{]}; no experiment distinguishing OI
from standard quantum theory, and no memory-witness constraint, is
established {[}Main §4.5{]}. For the cosmological constant: a derivation
of the covariant field equations with their source coupling. For
black-hole information: whether H-scramble and H-decouple hold, the Page
value on a proper energy shell and along a single trajectory, and, in
{[}GR{]}'s ladder, the covariant field equations.

Further open problems are: (1) machine checks of the steps of §2.2, §2.4
and §2.5 that lack one; (2) whether the minimal hidden predictive memory
equals \(M_t\) once ledger bits that cannot influence the future are
quotiented (the fixed-dimension form is false) {[}Main §3.4{]}; (3)
exact gluing for recurrent target families, and approximation on windows
short compared with the recurrence time {[}Main §3.4{]}; (4) (C2)
necessity, which is conditional on a hidden-sector mixing hypothesis
motivated by eigenstate thermalization {[}Main §3.4{]}.

\hypertarget{conclusion}{%
\subsubsection{6.4 Conclusion}\label{conclusion}}

At every finite horizon the laws an embedded observer of a finite
reversible system can record are exactly the finite stochastic laws, and
exactly the fixed-basis unitary Born laws; because the representation is
universal, the discriminating content lies in memory, which every
deterministic completion carries in its hidden state. Exact operational
quantum mechanics is reached only under completion conditions that OI
alone does not imply, and representability does not select the evolution
relating two times.

Applied to vacuum energy, OI proposes that the zero-point sum of the
emergent description is not a gravitational source, conditional on
covariant field equations that are not derived; applied to black-hole
information, it separates conservation from accessible recovery, while
the framework's Page curve is not established. Neither application is a
demonstrated resolution of its problem.

\hypertarget{code-data-and-formal-verification}{%
\subsection{Code, data and formal
verification}\label{code-data-and-formal-verification}}

The declarations in Appendix A were checked in Lean 4.33.0 with Mathlib
v4.33.0 at commit 17b830bcd0ecc36e10eff5b9eeca03bdd7198d86, whose Lean
sources are identical at the corpus commit \texttt{b64919cd}, by
building the Lean project in \texttt{verification/lean-mathlib}
(\texttt{lake\ exe\ cache\ get}, \texttt{lake\ -\/-rehash\ build}) and
running \texttt{python3\ tools/release\_gate.py}, whose
\texttt{lean-axioms} step requires every \texttt{\#print\ axioms} report
to lie within \texttt{propext}, \texttt{Classical.choice} and
\texttt{Quot.sound}. Each declaration in Appendix A reports exactly
these three; the expected result and the comparison recipe are in
\texttt{papers/introduction\_repro/}.

The probes and the black-hole computations {[}BHa; BHb{]} are released
under the MIT License, at these repository paths at the corpus commit:

\begin{itemize}
\tightlist
\item
  \texttt{papers/oi\_lattice\_code/foundations/} (probes; also in the
  release archive, concept DOI 10.5281/zenodo.19060318);
\item
  \texttt{verification/programmes/black-hole-information/} (black-hole
  computations; not part of the release archive).
\end{itemize}

\hypertarget{appendix-a.-machine-checked-declarations}{%
\subsection{Appendix A. Machine-checked
declarations}\label{appendix-a.-machine-checked-declarations}}

Modules of the \texttt{OIBridge} Lean library; line numbers at commit
\texttt{17b830bc}.

\begingroup\small\raggedright

\begin{itemize}
\tightlist
\item
  \textbf{Theorem 1} (§2.2): \texttt{finite\_horizon\_equivalence}
  (Equivalence:462); directions \texttt{S\_imp\_D} (:422),
  \texttt{D\_imp\_Qfb} (:294), \texttt{Qfb\_imp\_S} (:245), compositions
  :466--468; supporting \texttt{permMatrix\_mem\_unitaryGroup}
  (EquivalenceChain:179). Prior independence is proved on paper.
\item
  \textbf{Theorem 2} (§2.4): \texttt{pushforward\_identity},
  \texttt{distinguishability\_floor}, \texttt{capacity\_floor},
  \texttt{unavoidable\_hidden\_predictive\_memory} (HiddenMemory:83,
  :114, :170, :185), for a joint law of history, present and hidden
  state; the time-\(t\) identification and the standard vanishing
  criterion for conditional mutual information {[}CT06{]} are not
  machine-checked.
\item
  \textbf{Capacity bound} (§2.4): \texttt{c3\_necessity},
  \texttt{card\_hidden\_ge\_two\_pow\_Istar} (C3Necessity:80, :117).
\item
  \textbf{Stochastic-inverse lemma; Theorem 3} (§2.5):
  \texttt{isPermMatrix\_of\_stochastic\_inverse} (EquivalenceChain:110);
  \texttt{pIndivisible\_of\_rooted\_identity\_return\_of\_overlap}
  (RecurrenceHorizon:148).
\item
  \textbf{Readback forces indivisibility} (§2.5):
  \texttt{routed\_forces\_indivisible\_somewhere}
  (PhysicalC4Discharge:122) for routed readback; on paper for (C4).
\item
  \textbf{Relating evolution} (§3.3):
  \texttt{gl2\_strong\_gauge\_moves\_relative\_candidate},
  \texttt{gl3\_constant\_gauge\_preserves\_relative}
  (CoherentLiftGauge:537, :451); \texttt{bridge\_traj},
  \texttt{rooted\_eq\_iff\_slice\_eq} (forward direction),
  \texttt{realData\_traj\_stochastic} (TrackBQfbBridge:276, :180, :286).
\item
  \textbf{Theorem 4} (§3.2): \texttt{exactAll\_iff\_physical\_general}
  (GeneralCarrier:100); \(\Rightarrow\) \texttt{physical\_of\_exactAll}
  (PhysicalCharacterization:301); \(\Leftarrow\)
  \texttt{exactAll\_of\_physical\_general} (GeneralCarrier:93) via
  \texttt{krausSoundExt\_of\_validity\_inert} (OperationalValidity:124),
  \texttt{compositeCompleteness\_unconditional} (IsometryExtension:179,
  premise discharged by \texttt{finiteIsometryExtensionSF\_discharged},
  :116) and \texttt{exactAll\_of\_levelOne} (LevelOneSeam:196).
\item
  \textbf{Independence; OI alone does not imply QM} (§3.2, qubit):
  \texttt{five\_way\_minimality} (RankGapTheory:702),
  \texttt{oi\_alone\_not\_qm} (GeneralCarrier:160).
\end{itemize}

\hypertarget{appendix-b.-examples}{%
\subsection{Appendix B. Examples}\label{appendix-b.-examples}}

\hypertarget{b.1-a-minimal-model}{%
\subsubsection{B.1 A minimal model}\label{b.1-a-minimal-model}}

Let \(\mathcal C_V=\{0,1\}\) and \(\mathcal C_H=\{1,\dots,6\}\) {[}Main
§2.4{]}. For \(h\in\{1,2\}\), \(\varphi\) swaps \((0,h)\) with
\((1,h)\). It flips the visible bit and keeps \(h\), which then records
the flip. At fixed \(x\), \(\varphi\) swaps \(h=3\) with \(4\), and
\(5\) with \(6\). A second application reads the record and flips the
bit back. With \(h\) uniform: \[
\Gamma_1=\begin{pmatrix}2/3&1/3\\1/3&2/3\end{pmatrix},\qquad \Gamma_2=I,\qquad
\Gamma_1^{-1}\Gamma_2=\begin{pmatrix}2&-1\\-1&2\end{pmatrix}.
\] The only candidate divisor has negative entries, so the family is
P-indivisible. A Markov chain with the same one-step matrix would
instead give \(\Gamma_1^2\), with diagonal entries 5/9 and off-diagonal
entries 4/9.

With \(x_0\) uniform and independent of \(h\),
\(P(x_2{=}0\mid x_1{=}1,x_0{=}0)=1\) while
\(P(x_2{=}0\mid x_1{=}1)=1/3\). For \(P(x_0{=}0)=a\), the second value
is \(a/(2-a)\).

\hypertarget{b.2-memory-without-short-window-indivisibility}{%
\subsubsection{B.2 Memory without short-window
indivisibility}\label{b.2-memory-without-short-window-indivisibility}}

Let \(X_0\) and \(X_1\) be independent uniform bits and
\(X_2=X_0\oplus X_1\). Every two-time transition matrix is
\(\tfrac12J\), so the family factors through \(\Lambda=I\) and is
P-divisible. Yet \(I(X_0;X_2\mid X_1)=1\) bit. Up to flipping the final
bit, this is the process of {[}RHP14, Eq. (15){]}; it appears
concurrently in {[}Dou26, App. A{]}.

\hypertarget{b.3-the-card-table}{%
\subsubsection{B.3 The card table}\label{b.3-the-card-table}}

The two-card game of {[}Book §1.11{]} uses two cards, each red or black:
a visible \emph{hole card} \(h\), and a face-down \emph{burn card}
\(b\), uniform and independent of \(h\). A \emph{street} is one shuffle
step, in which the two cards trade places: \(\varphi(h,b)=(b,h)\). A
\emph{peek} reveals the hole card and overwrites the burn card, which
holds the record, with a fresh card drawn uniformly and independently.

The probability that the hole card at showdown is the hole card as dealt
is obtained by exact enumeration of the equally likely cases:

\begin{longtable}[]{@{}ll@{}}
\toprule\noalign{}
Protocol & Probability \\
\midrule\noalign{}
\endhead
\bottomrule\noalign{}
\endlastfoot
one shuffle & 1/2 \\
two shuffles & 1 \\
two shuffles, with a peek between them & 1/2 \\
Markov prediction for two shuffles, \(\Gamma_1^2\) & 1/2 \\
\end{longtable}

Here \(\Gamma_1=\tfrac12J\) and \(\Gamma_2=I\). The two rows of
\(\Gamma_1\) are equal, so for every \(\Lambda\) the two rows of
\(\Gamma_1\Lambda\) are equal and \(\Gamma_1\Lambda\neq I\): the family
is P-indivisible.

The peek row depends on the rule: if the peek only reveals or replaces
the hole card, the probability is 1, so the 1/2 comes from overwriting
the record, not from looking. With anticorrelated cards (\(b=1-h\)),
\(\Gamma_1\) is a permutation and the game is divisible.

On card inputs the game is a classical stochastic process, and its
P-indivisibility is notion (2) of §3.1. Embed the game as a permutation
unitary acting on an arbitrary hole-card state \(\rho\), with the burn
card in the uniform diagonal state {[}Main §3.2{]}. The one-step channel
is then the constant channel
\(\rho\mapsto\operatorname{tr}(\rho)\,I/2\), which is
entanglement-breaking {[}Main §3.2, Remark (failure of (C1)){]}. The
two-step map is the identity channel. The family is therefore
CP-indivisible {[}Main §3.2, Theorem (CP-indivisibility, framework
scope){]}.

Following Bäcker, Beyer and Strunz {[}BBS24{]}, a two-step process has
classical memory when the second map can be written as an instrument
realizing the first map, followed by channels conditioned on its
outcome. With the hidden slot initialized to a fixed card instead, the
swap has the maps \(\rho\mapsto|0\rangle\langle0|\) and
\(\rho\mapsto\rho\) and admits no such decomposition {[}Main §4.5{]}.
With the uniform burn card it does: the instrument
\(\rho\mapsto\tfrac14\sigma_o\rho\sigma_o\), \(o\in\{0,x,y,z\}\), sums
to the one-step map by the identity
\(\sum_o\sigma_oX\sigma_o=2\operatorname{tr}(X)\,I\), and the channels
\(\rho\mapsto\sigma_o\rho\sigma_o\), conditioned on \(o\), return
\(\rho\).

A card arrangement violates no Bell inequality {[}Book §1.11; Main
§3.3{]}.

\hypertarget{references}{%
\subsection{References}\label{references}}

\textbf{Corpus} (repository commit \texttt{b64919cd}, concept DOI
10.5281/zenodo.19060318)

{[}Main{]} A. Maybaum, ``The Incompleteness of Observation: The
Equivalence of Quantum Mechanics and Embedded Observation,''
\texttt{papers/Main.md}.

{[}GR{]} A. Maybaum, ``ℏ, the Bekenstein-Hawking Entropy, and the
Running-Vacuum Form of Dark Energy from the Cosmological Horizon,''
\texttt{papers/GR.md}.

{[}Book{]} A. Maybaum, \emph{The Incompleteness of Observation},
\texttt{book/}.

{[}BHa{]} A. Maybaum, recovery from accessible records in finite
reversible models, unpublished research record, internally reviewed;
\texttt{bh01\_access\_v1/} in the black-hole computations' directory.

{[}BHb{]} A. Maybaum, ``Permutation dynamics versus quantum scrambling;
Page-time fraction,'' unpublished research record, internally reviewed;
\texttt{bh02\_scramble\_v1/} in the black-hole computations' directory.

\textbf{External}

{[}Bar23{]} J. A. Barandes, ``The Stochastic-Quantum Theorem,''
arXiv:2309.03085v2 (2023).

{[}Bar25{]} J. A. Barandes, ``The Stochastic-Quantum Correspondence,''
\emph{Philosophy of Physics} \textbf{3}, 8 (2025); arXiv:2302.10778v3.

{[}Bar26{]} J. A. Barandes, ``The Born Representation Theorem and the
Unistochastic Theorem,'' arXiv:2608.04354v1 (2026).

{[}BB12{]} J. C. Baez and J. Biamonte, ``Quantum techniques for
stochastic mechanics,'' arXiv:1209.3632v5 (2012).

{[}BBS24{]} C. Bäcker, K. Beyer, and W. T. Strunz, ``Local disclosure of
quantum memory in non-Markovian dynamics,'' \emph{Phys. Rev.~Lett.}
\textbf{132}, 060402 (2024); arXiv:2310.01205v2.

{[}BD16{]} F. Buscemi and N. Datta, ``Equivalence between divisibility
and monotonic decrease of information in classical and quantum
stochastic processes,'' \emph{Phys. Rev.~A} \textbf{93}, 012101 (2016);
arXiv:1408.7062v3.

{[}Bel64{]} J. S. Bell, ``On the Einstein Podolsky Rosen paradox,''
\emph{Physics Physique Fizika} \textbf{1}, 195 (1964).

{[}BJA17{]} B. Bylicka, M. Johansson, and A. Acín, ``Constructive method
for detecting the information backflow of non-Markovian dynamics,''
\emph{Phys. Rev.~Lett.} \textbf{118}, 120501 (2017); arXiv:1603.04288v2.

{[}BLP09{]} H.-P. Breuer, E.-M. Laine, and J. Piilo, ``Measure for the
degree of non-Markovian behavior of quantum processes in open systems,''
\emph{Phys. Rev.~Lett.} \textbf{103}, 210401 (2009); arXiv:0908.0238v2.

{[}Cal26{]} S. Calvo, ``On the stochastic-quantum correspondence,''
arXiv:2601.18720v2 (2026).

{[}CHSH69{]} J. F. Clauser, M. A. Horne, A. Shimony, and R. A. Holt,
``Proposed experiment to test local hidden-variable theories,''
\emph{Phys. Rev.~Lett.} \textbf{23}, 880 (1969).

{[}CT06{]} T. M. Cover and J. A. Thomas, \emph{Elements of Information
Theory}, 2nd ed.~(Wiley, 2006).

{[}dMU21{]} L. de Moura and S. Ullrich, ``The Lean 4 theorem prover and
programming language,'' in \emph{Proc. CADE 28} (Springer, 2021),
p.~625.

{[}Dou26{]} J. Doukas, ``On the emergence of quantum mechanics from
stochastic processes,'' arXiv:2602.22095v3 (2026).

{[}EGGH26{]} G. Egri, M. Gömöri, B. Gyenis, and G. Hofer-Szabó,
``Trajectory of probabilities, probability on trajectories, and the
stochastic-quantum correspondence,'' arXiv:2602.23491v1 (2026).

{[}EvEMU11{]} G. F. R. Ellis, H. van Elst, J. Murugan, and J.-P. Uzan,
``On the trace-free Einstein equations as a viable alternative to
general relativity,'' \emph{Class. Quantum Grav.} \textbf{28}, 225007
(2011); arXiv:1008.1196v2.

{[}GH03{]} E. F. Galvão and L. Hardy, ``Substituting a qubit for an
arbitrarily large number of classical bits,'' \emph{Phys. Rev.~Lett.}
\textbf{90}, 087902 (2003); arXiv:quant-ph/0110166v3.

{[}Har04{]} L. Hardy, ``Quantum ontological excess baggage,''
\emph{Stud. Hist. Philos. Mod. Phys.} \textbf{35}, 267 (2004).

{[}Haw76{]} S. W. Hawking, ``Breakdown of predictability in
gravitational collapse,'' \emph{Phys. Rev.~D} \textbf{14}, 2460 (1976).

{[}HKLP24{]} T. Heinosaari, O. Kerppo, L. Leppäjärvi, and M. Plávala,
``Simple information processing tasks with unbounded quantum
advantage,'' \emph{Phys. Rev.~A} \textbf{109}, 032627 (2024);
arXiv:2308.07727v4.

{[}KP14{]} N. Kaloper and A. Padilla, ``Sequestering the Standard Model
vacuum energy,'' \emph{Phys. Rev.~Lett.} \textbf{112}, 091304 (2014);
arXiv:1309.6562v2.

{[}LP17{]} D. A. Levin and Y. Peres, with contributions by E. L. Wilmer,
\emph{Markov Chains and Mixing Times}, 2nd ed.~(American Mathematical
Society, 2017).

{[}Mathlib20{]} The mathlib Community, ``The Lean mathematical
library,'' in \emph{Proc. ACM SIGPLAN Conf. Certified Programs and
Proofs} (ACM, 2020), p.~367.

{[}Milz20{]} S. Milz, D. Egloff, P. Taranto, T. Theurer, M. B. Plenio,
A. Smirne, and S. F. Huelga, ``When is a non-Markovian quantum process
classical?,'' \emph{Phys. Rev.~X} \textbf{10}, 041049 (2020);
arXiv:1907.05807v3.

{[}Mon08{]} A. Montina, ``Exponential complexity and ontological
theories of quantum mechanics,'' \emph{Phys. Rev.~A} \textbf{77}, 022104
(2008); arXiv:0711.4770v2.

{[}Pag93{]} D. N. Page, ``Average entropy of a subsystem,'' \emph{Phys.
Rev.~Lett.} \textbf{71}, 1291 (1993); arXiv:gr-qc/9305007v2.

{[}Pag13{]} D. N. Page, ``Time dependence of Hawking radiation
entropy,'' \emph{J. Cosmol. Astropart. Phys.} \textbf{09}, 028 (2013);
arXiv:1301.4995v3.

{[}Pen20{]} G. Penington, ``Entanglement wedge reconstruction and the
information paradox,'' \emph{J. High Energy Phys.} \textbf{09}, 002
(2020); arXiv:1905.08255v3.

{[}RHP14{]} Á. Rivas, S. F. Huelga, and M. B. Plenio, ``Quantum
non-Markovianity: characterization, quantification and detection,''
\emph{Rep.~Prog. Phys.} \textbf{77}, 094001 (2014); arXiv:1405.0303v2.

{[}Sch21{]} H.-J. Schmidt, ``Dilation of stochastic matrices by coarse
graining,'' arXiv:2106.03513v1 (2021).

{[}SGD19{]} P. Strasberg and M. García Díaz, ``Classical quantum
stochastic processes,'' \emph{Phys. Rev.~A} \textbf{100}, 022120 (2019);
arXiv:1905.03018v2.

{[}Sha79{]} A. Shamir, ``How to share a secret,'' \emph{Commun. ACM}
\textbf{22}, 612 (1979).

{[}tH14{]} G. 't Hooft, ``The cellular automaton interpretation of
quantum mechanics,'' arXiv:1405.1548v3 (2014).

{[}Vol06{]} G. E. Volovik, ``Vacuum energy: myths and reality,''
\emph{Int. J. Mod. Phys. D} \textbf{15}, 1987 (2006);
arXiv:gr-qc/0604062v4.

{[}Wei89{]} S. Weinberg, ``The cosmological constant problem,''
\emph{Rev.~Mod. Phys.} \textbf{61}, 1 (1989).

\endgroup

\end{document}
% source-sha256: beb0a5202419a71a8a633f5999f7f48a5e8768a922b7681b7576c470b936feb8
